%% This is file `elsarticle-template-1-num.tex',
%%
%% Copyright 2009 Elsevier Ltd
%%
%% This file is part of the 'Elsarticle Bundle'.
%% ---------------------------------------------
%%
%% It may be distributed under the conditions of the LaTeX Project Public
%% License, either version 1.2 of this license or (at your option) any
%% later version.  The latest version of this license is in
%%    http://www.latex-project.org/lppl.txt
%% and version 1.2 or later is part of all distributions of LaTeX
%% version 1999/12/01 or later.
%%
%% Template article for Elsevier's document class `elsarticle'
%% with numbered style bibliographic references
%%
%% $Id: elsarticle-template-1-num.tex 149 2009-10-08 05:01:15Z rishi $
%% $URL: http://lenova.river-valley.com/svn/elsbst/trunk/elsarticle-template-1-num.tex $
%%
\documentclass[review,12pt]{elsarticle}

%% Use the option review to obtain double line spacing
%% \documentclass[preprint,review,12pt]{elsarticle}

%% Use the options 1p,twocolumn; 3p; 3p,twocolumn; 5p; or 5p,twocolumn
%% for a journal layout:
%% \documentclass[final,1p,times]{elsarticle}
%% \documentclass[final,1p,times,twocolumn]{elsarticle}
%% \documentclass[final,3p,times]{elsarticle}
%% \documentclass[final,3p,times,twocolumn]{elsarticle}
%% \documentclass[final,5p,times]{elsarticle}
%% \documentclass[final,5p,times,twocolumn]{elsarticle}

%% The graphicx package provides the includegraphics command.
\usepackage{graphicx}
%% The amssymb package provides various useful mathematical symbols
\usepackage{amssymb}
%% The amsthm package provides extended theorem environments
%% \usepackage{amsthm}

%% The lineno packages adds line numbers. Start line numbering with
%% \begin{linenumbers}, end it with \end{linenumbers}. Or switch it on
%% for the whole article with \linenumbers after \end{frontmatter}.
\usepackage{lineno}

\usepackage{blindtext}
\usepackage[hyphens]{url}
\usepackage{hyperref}
  \hypersetup{
  bookmarks=true,         % show bookmarks bar?
  unicode=false,          % non-Latin characters in Acrobat bookmarks
  pdftoolbar=true,        % show Acrobat toolbar?
  pdfmenubar=true,        % show Acrobat menu?
  pdffitwindow=false,     % window fit to page when opened
  pdfstartview={FitH},    % fits the width of the page to the window
  pdftitle={Ludurel Service Agreement},    % title
  pdfauthor={sb swae},     % author
  pdfsubject={},   % subject of the document
  pdfcreator={sb swae},   % creator of the document
  pdfproducer={sb swae}, % producer of the document
  pdfkeywords={}, % list of keywords
  pdfnewwindow=true,      % links in new window
  colorlinks=true,        % false: boxed links; true: colored links
  linkcolor=Gray,         % color of internal links
  citecolor=Gray,        % color of links to bibliography
  filecolor=magenta,      % color of file links
  urlcolor=Gray           % color of external links
  }

%% natbib.sty is loaded by default. However, natbib options can be
%% provided with \biboptions{...} command. Following options are
%% valid:

%%   round  -  round parentheses are used (default)
%%   square -  square brackets are used   [option]
%%   curly  -  curly braces are used      {option}
%%   angle  -  angle brackets are used    <option>
%%   semicolon  -  multiple citations separated by semi-colon
%%   colon  - same as semicolon, an earlier confusion
%%   comma  -  separated by comma
%%   numbers-  selects numerical citations
%%   super  -  numerical citations as superscripts
%%   sort   -  sorts multiple citations according to order in ref. list
%%   sort&compress   -  like sort, but also compresses numerical citations
%%   compress - compresses without sorting
%%
%% \biboptions{comma,round}

% \biboptions{}

\journal{Journal Name}

\begin{document}

\begin{frontmatter}

%% Title, authors and addresses

\title{Management of Liability, Money, Logistics, and User Experience by Turo, Inc.}

%% use the tnoteref command within \title for footnotes;
%% use the tnotetext command for the associated footnote;
%% use the fnref command within \author or \address for footnotes;
%% use the fntext command for the associated footnote;
%% use the corref command within \author for corresponding author footnotes;
%% use the cortext command for the associated footnote;
%% use the ead command for the email address,
%% and the form \ead[url] for the home page:
%%
%% \title{Title\tnoteref{label1}}
%% \tnotetext[label1]{}
%% \author{Name\corref{cor1}\fnref{label2}}
%% \ead{email address}
%% \ead[url]{home page}
%% \fntext[label2]{}
%% \cortext[cor1]{}
%% \address{Address\fnref{label3}}
%% \fntext[label3]{}


%% use optional labels to link authors explicitly to addresses:
%% \author[label1,label2]{<author name>}
%% \address[label1]{<address>}
%% \address[label2]{<address>}

\author{sb swae}

\address{\href{mailto:sb.swae@gmail.com}{sb.swae@gmail.com}; \url{https://swaevior.io}}

\begin{abstract}
%% Text of abstract
\blindtext
\end{abstract}

\begin{keyword}
Turo \sep car sharing \sep peer-to-peer \sep liability \sep insurance \sep money \sep user experience \sep user interface
%% keywords here, in the form: keyword \sep keyword

%% MSC codes here, in the form: \MSC code \sep code
%% or \MSC[2008] code \sep code (2000 is the default)

\end{keyword}

\end{frontmatter}

%%
%% Start line numbering here if you want
%%
\linenumbers

%% main text
\section{Introduction and Background}
\label{S:1}

\blindtext

\begin{itemize}
\item Bullet point one
\item Bullet point two
\end{itemize}

\begin{enumerate}
\item Numbered list item one
\item Numbered list item two
\end{enumerate}

\subsection{Subsection One}

\blindtext

\begin{table}[h]
\centering
\begin{tabular}{l l l}
\hline
\textbf{Treatments} & \textbf{Response 1} & \textbf{Response 2}\\
\hline
Treatment 1 & 0.0003262 & 0.562 \\
Treatment 2 & 0.0015681 & 0.910 \\
Treatment 3 & 0.0009271 & 0.296 \\
\hline
\end{tabular}
\caption{Table caption}
\end{table}

\subsection{Subsection Two}

\blindtext

\begin{figure}[h]
\centering\includegraphics[width=0.4\linewidth]{placeholder}
\caption{Figure caption}
\end{figure}

\blindtext

\begin{equation}
\label{eq:emc}
e = mc^2
\end{equation}

\section{Liability Exposure and Insurance Coverage}
\label{S:2}
\subsection{Owner's Insurance}
Turo offers three protection plans for owners who make their cars available for rent: basic, standard, and premium. The owner's choice of protection plan determines the share of the rental fee that the owner receives. Owners are also permitted to simply carry their own insurance on their vehicle.

Insurance coverage in the United States is provided by a group plan administered by Liberty Mutual.

Turo openly admits that their protection plans come with an inherent level of uncertainty. Turo senior claims manager Chris Aragon states that, ``If there’s an engine failure, and it’s something that’s caused by a mechanical failure and not something that the renter could have caused by using the vehicle, that’s something that is not covered by us. That’s something that’s just a mechanical breakdown that you’d be expected to pay for.''
  \footnote{
  Kristen Lee. Jalopnik, 22 February 2017. \url{https://jalopnik.com/how-insurance-works-when-you-rent-out-your-car-on-turo-1792401490}
  }
This is consistent with standard insurance practices in the United States. Insurance carriers do not typically provide coverage for mechanical failures in vehicles as these cases are often covered by warranties. In the event that a component of the vehicle fails due to driver abuse, Turo states that it will determine the cause and act accordingly.

\subsubsection{Basic Protection}
The basic protection plan offers \$1,000,000 USD in liability insurance, covers physical damage to the vehicle up to a \$125,000 USD, and has a \$3,000 USD deductible. Private auto insurance deductibles in the United Staes typically range between \$100 and \$1000 USD, though they can be as high as \$2500 in some cases. Under this plan, Turo will pay 20\% up to the first \$3,750 of a damage claim and then covers 100\% beyond that limit, up to a cap of \$125,000. At this level, owners receive 85\% of the trip fee.

\subsubsection{Standard Protection}
The standard protection plan offers \$1,000,000 USD in liability insurance, covers physical damage to the vehicle up to \$125,000, and has a \$0 deductible. Owners also receive reimbursement for a replacement vehicle during the period of time that their personal vehicle is being repaired. At this level, owners receive 75\% of the trip fee.

\subsubsection{Premium Protection}\label{owner-premium-protection}
The premium protection plan offers \$1,000,000 USD in liability insurance, covers physical damage to the vehicle up to \$125,000, and has a \$0 deductible. Owners also receive reimbursement for a replacement vehicle during the period of time that their personal vehicle is being repaired, includes coverage for wear and tear on the exterior of the vehicle, and also includes coverage for the loss of rental income.
\footnote{Turo determines ``lost rental income'' by taking the mean daily rental earnings of the owner for the past 60 days and multiplying that dollar amount by the number of days that the owner's vehicle spent being repaired within a ``reasonable range.'' Turo's ``reasonable range'' is determined like so: Turo assumes that a mechanic should reasonably be able to put in four hours of labor per day to repair the vehicle. Turo takes the number of labor hours on the mechanic's estimate and divides it by four. This ``number of days'' is then multiplied by the 60-day mean rental earnings. In the event that the owner cannot pick up otherwise obtain the vehicle within the calculated number of days, Turo will not pay for futher days of lost rental income.
}
At this level, owners receive 65\% of the trip fee.

\subsubsection{Owner-provided Protection}
Vehicle owners also have the option of opting-out of any coverage provided that they carry commercial insurance on their own. Such an insurance policy is typically only carried by businesses. At this level, owners receive 90\% of the trip fee and have no protection whatsoever provided to them by Turo or Liberty Mutual.

\subsection{Renter's Insurance}
Renters are presented with a trio of insurance coverage options.

Physical damage protection includes both comprehensive (property damage or loss to a vehicle, including theft, fire, impacts with wildlife,  vandalism, and acts of nature) and collision (protection in the event that two or more vehicles impact each other or in the event that a single vehicle suffers and impact with an object) coverage.

Personal injury protection varies widely by state, but typically provides coverage for medical expenses, lost income, funeral expenses, and the like for for the driver and passenger in the vehicle with no regard given to who is at fault.

Uninsured/underinsured motorist coverage pays the renter of the renter's passengers for bodily injury they may suffer in the event that they experience a collision caused by a driver without insurance, a hit-and-run driver, a driver who has exhausted the benefits available to them under their own insurance plan, or a driver whose insurance carrier is found to be insolvent and therefore incapable of paying a claim.

\subsubsection{Premium Rental Insurance}\label{premium-rental-insurance}
The premium rental insurance option provides \$1,000,000 USD in bodily injury and property damage coverage to renters. Personal injury protection coverage (also known as no-fault or first party benefits) are capped at the statutory minimum required in the state where the vehicle is registered. In some cases, this amount may be \$0 USD. Uninsured/underinsured motorist coverage is also capped at the statutory minimum required in the state where the vehicle is registered. In some cases, this may be \$0 USD. The renter's out-of-pocket exposure is limited to a deductible payment of \$500, which is collected when the owner reports damage. In the event that the total costs are less than \$500, then the renter will be refunded the difference.

\subsubsection{Basic Rental Insurance}\label{basic-rental-insurance}
The basic rental insurance option provides \$1,000,000 USD in bodily injury and property damage coverage to renters. Personal injury protection coverage (also known as no-fault or first party benefits) are capped at the statutory minimum required in the state where the vehicle is registered. In some cases, this amount may be \$0 USD. Uninsured/underinsured motorist coverage is also capped at the statutory minimum required in the state where the vehicle is registered. In some cases, this may be \$0 USD. The renter's out-of-pocket exposure is limited to a deductible payment of \$3000 USD. In the event that the total costs are less than \$500, then the renter will be refunded the difference.

\subsubsection{Declining Rental Insurance Coverage}\label{commercial-insurance}
Renters are also allowed to decline a protection plan. In this case, the renter's exposure is effectively unlimited and is only bounded by the actual cash value of the vehicle plus all related costs, as determined by the process described in \S\ref{claims-process}. In the event that a rented vehicle is lost, stolen, or deemed a total loss,
\footnote{For Turo's purpose, a ``total loss'' is defined by anticipated repair costs in excess of 75\% of the vehicle's actual cash value.
}
then the renter is liable for paying the entire actual cash value of the vehicle, plus all related costs and minus any residual salvage value.
\footnote{Salvage value is the expected sale value of an asset at the end of its useful life. Insurance companies may employ a number of methods to determine the salvage value of a vehicle, but it is often classified as the sum value of the remaining usable parts in a vehicle.
}

\subsubsection{Other Sources of Renter's Insurance}
Some renters may be able to obtain, or already have as a part of their personal auto insurance policy, coverage that extends to vehicles rented by the primary policy holder.

Some renters may be also be able to obtain, or already have as a part of their existing credit card account, coverage that extends to vehicles rented by the cardholder using that credit card as the payment method.

Such coverage varies widely between insurance carriers and credit providers and may, in some cases, be deficient. Turo states that it is ``highly unlikely that [the renter] would be covered by any credit card insurance when [they] book cars through the Turo marketplace.''
\footnote{\url{https://support.turo.com/hc/en-us/articles/203990610-I-d-like-a-detailed-explanation-of-insurance-and-protection-provisions}}

\subsubsection{Subrogation}
Insurance offered by Turo is secondary to any existing coverage that the renter may have through other sources (\emph{i.e.,} that available through a credit card provider, personal auto policy, or other insurance policy). Practically speaking, the secondary nature of Turo's rental insurance offerings means that the renter is primarily liable, but can satisfy their obligations through other sources. In the event of a claim, Turo will first attempt to collect deductibles directly from the renter and then then seek reimbursement from other sources of coverage that the rent may have. (In US insurance practice, this process is known as subrogation.
\footnote{Subrogation is the legal doctrine through which one party is entitled to enforce the legal rights or protections of another party for its own benefit.
}
)

\subsubsection{Exception to Renter's Elected Protection Plans}
In the event that the renter is found to have violated Turo's terms of service, engaged in prohibited uses of the rented vehicle, or is otherwise found to have recklessly used the rented vehicle, the deductible limits mentioned in \S\ref{premium-rental-insurance} and \S\ref{basic-rental-insurance} do not apply and the renter will be liable for physical damage up the full actual cash value of the vehicle including Turo's related costs.

\subsection{Claims Process}\label{claims-process}
A renter is supposed to report any incidents involving a rented vehicle to the relevant authorities. Urgent matters should be referred to the local police and either to Turo or to the owner of the vehicle within 24 hours of the incident.

In the event that a host does not learn about an incident that occurred with their vehicle during a trip until after the vehicle has been returned, the owner has 24 hours to report the incident to Turo in order to remain eligible for Turo's coverage.

If a vehicle's owner has elected to decline Turo's coverage, then they are responsible for making a claim directly to their own insurance company.

If the owner of a vehicle has elected either the premium or the standard Turo protection plans, then they are eligible to receive \$30 USD per day for up to 10 days (\$300 USD cap) to rent a replacement vehicle. They may also elect to receive a travel credit from Turo that may be used to rent a vehicle on the Turo marketplace. Additionally, users may submit receipts from public transport, taxis, and ride services like Uber or Lyft for reimbursement.

Owners and renters have the option to resolve claims directly with each other. However, if an owner elects to file a claim through Turo, they have until 24 hours after the end of a trip to visit \url{https://turo.com/resolutions} to file an eligible claim. Turo states that one of their claims specialists will make contact with the owner of the vehicle within 24 of a claim being submitted. Turo will then provide instructions to the parties about how to obtain an appraisal.

Turo uses a third-party administrator called the Littleton Group to determine a vehicle's actual cash value. Littleton, in turn, uses AutoClaims Direct, a nation-wide network of independent auto appraisers. In the event of a claim, AutoClaims Direct will dispatch an appraiser to perform a visual inspection of the vehicle and to obtain photographic documentation of the vehicle's condition. The appraiser will also collect information on the general condition of the vehicle, the odometer reading, and any other information that may help accurately determine the value of the vehicle.

Software is employed to sweep across the market to obtain information on comparable vehicles to the one subject to a claim. The software will then determine the value of the vehicle based on how it compares in condition and mileage to comparable vehicles on the market at that time. It may also include a factor that accounts for depreciation.

If a claim made to Turo is determined to be eligible, Turo notifies the renter and charges their payment method(s) for an initial claim processing cost of up to \$575.

The owner is then presented with their options for resolution. They may elect to not pursue, the owner and the renter can elect to resolve the issue directly with each other, the owner may resolve the claim directly with the insurance carrier of the renter (or a third party in the event that the driver is found no to be at fault), or the claim may be processed through Turo's claims administrator.

Depending on the resolution option elected by the parties, the owner's car may be repaired and the renter may be responsible for settling their financial obligations, including payment of their deductible.

\section{Data and Metrics}
How is the control of the distance traveled (if I travel more than the established number of miles, how will the service know about it and how money is written off; Is there a computer/logger/tracker inside the vehicle?)
Can the car owner track the location of the car in real time? Alexander is confident that this is possible.
Is there a gasoline and mileage reader in the car?

\section{Logistics}
How is the meeting with the owner (at what time are the contact data transmitted, what does the owner say at the meeting? How fast is this information exchanged? - what are the recommendations)
How to set where you need to leave (return) the vehicle. Is it possible to drop the car off in another location?
Is the renter responsible for filling the car up before returning it? Is the renter billed for the fuel cost? How does Turo establish how much fuel is needed?

\section{Money Management}
At what point in time is the money is debited from the renter? (immediately or after the return of the machine)
Is there a security deposit taken at the time of reservation? How quickly is the deposit returned? Or do they reserve the right to charge your card at any time? What if the card Tx fails?

\section{User Experience}

If the Tenant did not like the car? If payment occurred before, how does the money refund happen? Canceling reservations?
Similarly, for the cleanliness of the car - if I rent a car in a dirty condition, how to find out about this service and how payment is built. How does Turo know this information Does owner provide this information to Turo? If so, how? ? How do they mediate this?

\subsection{User Responsibilities and Obligations}
  \subsubsection{Wear and Tear}
    \paragraph{Owners}
    ``Normal wear and tear'' is the result of the normal operation of a vehicle and therefore should be an expected part of the experience of sharing or renting a car on Turo. Owners are not, under any circumstances, protected against normal wear and tear of the interior of their vehicle(s). Owners who have elected the premium protection plan (\S\ref{owner-premium-protection}) receive protection against normal exterior wear and tear.
      \footnote{Defined by Turo as ``any dings, dents, cracks, or scratches to the exterior body of the vehicle that is 3 inches in diameter or less. This includes, but is not limited to, rims, wheels, hubcaps, any painted or textured area for the body of the vehicle, and moldings.'' See \url{https://support.turo.com/hc/en-us/articles/217043898}.
      }
    \paragraph{Renters}
    Renters are not responsible for wear and tear that are the result of normal use of a vehicle.

    Renters are fully responsible for damage that is the result of misuse or prohibited use of a vehicle, ``significant damage'' to the interior of a vehicle, and mechanical damage due to negligence, intentional acts, or improper driving on the part of the renter.
      \footnote{Owners who decline any of Turo's in-house protection plans and instead opt to provide their own commercial insurance (see \S\ref{commercial-insurance}) are not covered by Turo against any of the aforementioned items.
      }

  \subsubsection{Mechanical Failure}
  If the machine itself broke down, what actions does Turo take to pick up the car. (Do they pick up the car? Are there any charges?)
  \subsection{Traffic Violations and Tickets}
  If the renter violates the traffic rules, who is fined or ticketed and how? (where fines come, how to pay it first and how then compensation takes place)
    \paragraph{Renters}
    Renters are responsible for paying all tickets for violations that they commit during their trip, with the exception of moving violations and photo tickets.
      \footnote{In the event that a vehicle is ticketed for speeding, running a red light, or other moving violations, including those captured by cameras, and the owner of the vehicle receives said ticket via mail, then Turo agrees that it will provide the information required by the ticketing body or to the owner such that the owner may contest the ticket and transfer liability for the moving violation to the renter.

      Renters are not supposed to pay moving violations or photo tickets immediately because the liability for such violations (including ``points'' assigned in many US states) may be assigned to the owner of the vehicle for a traffic offense that they did not actually commit.
      }
    Renters are also responsible for parking tickets, towing fees, and fines that may be levied for up to 24 hours after the end of a trip, provided that such tickets, fines, and/or fees are a result of improper parking on the behalf of the renter.

    The renter is responsible for informing the owner in the event that they receive a ticket during their trip and ought to immediately pay it.

    \paragraph{Owners}
    In the event that an owner notifies Turo that a ticket was received during the renter's trip or within the 24 hour period thereafter, Turo will charge the renter's payment method(s) for the total cost of the ticket plus administration fees levied by Turo. Owners must submit their requests for reimbursement within 90 days after the trip in question ended, provide documentation of the ticket, and receipts for payment of the ticket.

\section{User Interface}
We need all screenshots of the renter’s private office (after the trip is completed). I need screenshot of every page inside the app (Mobile and desktop)
If it works: personal car owner's office - what he sees there, screenshots of the pages.
Is there a single app for all users or do owners and renters each have dedicated apps?

  \subsection{Renters}
  \subsection{Owners}




%% The Appendices part is started with the command \appendix;
%% appendix sections are then done as normal sections
%% \appendix

%% \section{}
%% \label{}

%% References
%%
%% Following citation commands can be used in the body text:
%% Usage of \cite is as follows:
%%   \cite{key}          ==>>  [#]
%%   \cite[chap. 2]{key} ==>>  [#, chap. 2]
%%   \citet{key}         ==>>  Author [#]

%% References with bibTeX database:

\bibliographystyle{model1-num-names}
\bibliography{sample.bib}

%% Authors are advised to submit their bibtex database files. They are
%% requested to list a bibtex style file in the manuscript if they do
%% not want to use model1-num-names.bst.

%% References without bibTeX database:

% \begin{thebibliography}{00}

%% \bibitem must have the following form:
%%   \bibitem{key}...
%%

% \bibitem{}

% \end{thebibliography}


\end{document}

%%
%% End of file `elsarticle-template-1-num.tex'.
